\documentclass{article}

\usepackage{amsmath,amssymb}
\usepackage{xurl}
\usepackage[hidelinks]{hyperref}

% Shared commands for notes in docs/paper-gaps/.

\DeclareUnicodeCharacter{2080}{\ifmmode{}_{0}\else\textsubscript{0}\fi}
\DeclareUnicodeCharacter{2081}{\ifmmode{}_{1}\else\textsubscript{1}\fi}
\DeclareUnicodeCharacter{2082}{\ifmmode{}_{2}\else\textsubscript{2}\fi}
\DeclareUnicodeCharacter{2099}{\ifmmode{}_{n}\else\textsubscript{n}\fi}

\newcommand{\Lean}{\textsc{Lean}}
\ExplSyntaxOn
\NewDocumentCommand{\leanid}{m}{
  \ifmmode
    \textsf{\detokenize{#1}}
  \else
    \group_begin:
    \urlstyle{sf}
    \tl_set:Nn \l_tmpa_tl {#1}
    \tl_replace_all:Nnn \l_tmpa_tl {\_} {_}
    \exp_args:NV \path \l_tmpa_tl
    \group_end:
  \fi
}
\ExplSyntaxOff
\providecommand{\tr}{\operatorname{tr}}
\newcommand{\tnleanIssueBase}{https://github.com/LionSR/TNLean/issues/}
\newcommand{\tnleanPrBase}{https://github.com/LionSR/TNLean/pull/}
\newcommand{\tnleanDocBase}{https://sirui-lu.com/TNLean/docs/}
\newcommand{\ghissue}[1]{\href{\tnleanIssueBase#1}{GitHub issue \##1}}
\newcommand{\ghpr}[1]{\href{\tnleanPrBase#1}{PR \##1}}
\newcommand{\doclink}[1]{\href{\tnleanDocBase#1.html}{\path{#1}}}

% Dirac notation and the identity operator, as in blueprint/src/macros/common.tex.
% \providecommand keeps note-local definitions authoritative.
\usepackage{dsfont}
\providecommand{\ket}[1]{|#1\rangle}
\providecommand{\bra}[1]{\langle #1|}
\providecommand{\braket}[2]{\langle #1 | #2 \rangle}
\providecommand{\ketbra}[2]{|#1\rangle\!\langle #2|}
\providecommand{\Id}{\mathds{1}}
\providecommand{\one}{\mathds{1}}
\providecommand{\C}{\mathbb{C}}
\providecommand{\MN}[1]{M_{#1}(\C)}
\providecommand{\MD}{M_D(\C)}

% External referencing for paper sources.  The build helper writes sanitized
% auxiliary files under build/paper-aux/ and excludes duplicated source labels.
\usepackage{xr}

% Import a generated paper auxiliary file with a caller-supplied label prefix.
% The candidate paths support builds from docs/paper-gaps/, the repository root,
% and build/paper-aux/ itself.
\newcommand{\importpaperaux}[2]{%
  \IfFileExists{../../build/paper-aux/#2.aux}{%
    \externaldocument[#1]{../../build/paper-aux/#2}%
  }{%
    \IfFileExists{build/paper-aux/#2.aux}{%
      \externaldocument[#1]{build/paper-aux/#2}%
    }{%
      \IfFileExists{#2.aux}{%
        \externaldocument[#1]{#2}%
      }{}%
    }%
  }%
}

% General external reference macros
\newcommand{\paperref}[2]{\ref{#1-#2}}
\newcommand{\papereqref}[2]{(\ref{#1-#2})}

% CPSV16 source (1606.00608 / MPDO-22-12-17-2.tex) external document and shortcuts
\importpaperaux{cpsv16-}{1606.00608/MPDO-22-12-17-2-tnlean}
\newcommand{\cpsvref}[1]{\paperref{cpsv16}{#1}}
\newcommand{\cpsveqref}[1]{\papereqref{cpsv16}{#1}}

% Verdict marker: \gapnote{<kind>}{<status>}. Kind is the policy
% classification (clarification, local-correction, scope-restriction,
% unfaithful, false-source, open-gap); status is open, wip (elimination
% underway), resolved, or historical. Machine-read by the site index;
% prints a verdict line.
\newcommand{\gapnote}[2]{\par\noindent\textbf{Verdict:} #1 (#2).\par}


\title{The Commuting-Form-to-SAL Implication for Simple MPDOs}
\author{}
\date{2026-07-31}

\begin{document}
\maketitle
\gapnote{scope-restriction}{resolved}

\section{Source statement}

The proposition labelled \emph{4to2} in Appendix~C.2, lines~1597--1619, of
Ref.~\cite{Cirac2016MPDO_arXiv} considers a normal MPDO block
$\mathcal K$. Its finite-chain operators are positive scalar multiples of
products of translates of one positive two-site bond:
\begin{align}
    \sigma^{(N)}(\mathcal K)
        &= c_N\prod_{i=0}^{N-1}B_{i,i+1},
    \qquad c_N>0,
    \label{eq:cpsv16_commuting_form_to_sal_source_bond_product}\\
    [B_{i,i+1},B_{j,j+1}]
        &=0.
    \label{eq:cpsv16_commuting_form_to_sal_source_bond_commutation}
\end{align}
Together with zero correlation length, this hypothesis is asserted to imply
saturation of the area law.

The phrase ``of the form'' refers to the single bond in the source equation
labelled \emph{expression}, at lines~1571--1578 of
Ref.~\cite{Cirac2016MPDO_arXiv}. It therefore specifies one
translation-invariant bond family, not a separate existence statement at each
chain length.

\section{The theorem cited through Beigi}

The reference denoted by \emph{Beigi} in
Ref.~\cite{Cirac2016MPDO_arXiv} is Salman Beigi,
\emph{Classification of the phases of 1D spin chains with commuting
Hamiltonians}, arXiv:1105.1019v2, J. Phys. A \textbf{45} (2012) 025306,
\href{https://doi.org/10.1088/1751-8113/45/2/025306}{%
doi:10.1088/1751-8113/45/2/025306}.

The relevant result is Lemma~2.1, labelled \emph{lem:comm} in that arXiv
source and credited there to Bravyi and Vyalyi. If Hermitian operators
$X_{AB}$ and $Y_{BC}$ commute after their natural embeddings, the middle
space admits a decomposition
\begin{align}
    \mathcal H_B
        &\cong
        \bigoplus_{\beta\in V}
        \mathcal H_{\beta_l}\otimes\mathcal H_{\beta_r},
    \label{eq:cpsv16_commuting_form_to_sal_middle_space_decomposition}
\end{align}
such that $X_{AB}$ acts trivially on every $\mathcal H_{\beta_r}$ and
$Y_{BC}$ acts trivially on every $\mathcal H_{\beta_l}$. The application in
Ref.~\cite{Cirac2016MPDO_arXiv} uses neighboring translates of one local
interaction. Translation invariance permits the same one-site decomposition
at every site, independently of the system size.

The proof uses the representation theory of finite-dimensional
$C^*$-algebras. It forms the commutant of the operator-Schmidt factors of
$X_{AB}$ and identifies that algebra, after a unitary change of basis, with
\begin{align}
    \bigoplus_{\beta\in V}
        \mathbf 1_{\beta_l}\otimes
        \mathbf L(\mathcal H_{\beta_r}).
    \label{eq:cpsv16_commuting_form_to_sal_middle_algebra_normal_form}
\end{align}
The commutant acts complementarily on the left tensor factors.

\section{Present formal distinction}

The predicate \leanid{MPOTensor.HasCommutingForm} is chainwise: for each
$N\geq2$, it permits a new two-site bond $B_N$. This predicate is weaker than
the source hypothesis and does not support the comparison of chains of lengths
$N$, $N+1$, and $N+2$ used in the source trace-factorization argument.

The formal model instead fixes one bond whose commuting translates realize
every finite chain.\footnote{The bond family is
\leanid{MPOTensor.TranslationInvariantBondData}; the realization equalities
are fields of \leanid{MPOTensor.EtaLocalStructureData}.}

The source ZCL hypothesis is \leanid{MPOTensor.IsSourceZCL}, the
scale-invariant idempotence condition for the physical-trace transfer. The
predicate \leanid{MPOTensor.IsZCL} concerns the doubled-index completely
positive transfer map and is not a faithful substitute for Proposition
\emph{4to2} of Ref.~\cite{Cirac2016MPDO_arXiv}.

\section{Resolution of the implication}

\subsection{Finite-dimensional spatial decomposition}

The finite-dimensional $C^*$-algebra component of the Bravyi--Vyalyi argument
quoted in Ref.~\cite{Cirac2016MPDO_arXiv} is formalized. For a
star-subalgebra $S\subseteq\mathbf L(H_B)$, one orthonormal identification
\begin{align}
    H_B
        &\cong
        \bigoplus_{\beta}
        H_{\beta_l}\otimes H_{\beta_r}
    \label{eq:cpsv16_commuting_form_to_sal_star_subalgebra_space_decomposition}
\end{align}
simultaneously places $S$ and its commutant $S'$ in
\begin{align}
    S
        &\subseteq
        \bigoplus_\beta
        \mathbf 1_{\beta_l}\otimes\mathbf L(H_{\beta_r}),
    \label{eq:cpsv16_commuting_form_to_sal_star_subalgebra_right_action}\\
    S'
        &\subseteq
        \bigoplus_\beta
        \mathbf L(H_{\beta_l})\otimes\mathbf 1_{\beta_r}.
    \label{eq:cpsv16_commuting_form_to_sal_star_commutant_left_action}
\end{align}
This complementary action is formalized by
\leanid{StarSubalgebra.exists_complementary_action_orthonormalBasis}.

The coefficient step in the cited proof is also available. Fixing the outer
indices produces middle-site matrices $X_{aa'}$ and $Y_{cc'}$. Commutation of
the overlapping operators gives
\begin{align}
    [X_{aa'},Y_{cc'}]
        &=0
        \qquad\text{for all }a,a',c,c'.
    \label{eq:cpsv16_commuting_form_to_sal_middle_coefficient_commutation}
\end{align}
Applying the complementary-action theorem to the star-commutant of the
$Y_{cc'}$ gives one orthonormal decomposition on which all $X_{aa'}$ and
$Y_{cc'}$ act on complementary factors. This result is
\leanid{Matrix.exists_middle_spatial_decomposition_of_overlappingLifts_commute}.

The conjugated coordinate form of Lemma~2.1 cited in
Ref.~\cite{Cirac2016MPDO_arXiv} is
\leanid{Matrix.exists_unitary_blockActions_of_overlappingLifts_commute}.
It constructs one unitary identification of the middle space and combines the
coefficient matrices into Hermitian operators $R_{A\beta_l}$ and
$S_{\beta_rC}$. Applying the inverse coordinate identifications and cancelling
the unitary factors gives the two equalities for the original operators. This
is
\leanid{Matrix.exists_unitary_blockSum_equalities_of_overlappingLifts_commute}.
Its hypotheses are precisely Hermiticity of both overlapping operators and
commutation of their natural three-factor lifts. The block operators are
Hermitian; no unitarity property is asserted for them.

The cyclic three-site coordinate identification is also available. The
equivalence \leanid{MPOTensor.threeSiteOverlappingEquiv} sends a three-site
configuration to the ordered left-associated triple. Under this common
equivalence,
\leanid{MPOTensor.EtaLocalStructureData.reindex_bondAt_zero_eq_leftOverlappingLift}
and
\leanid{MPOTensor.EtaLocalStructureData.reindex_bondAt_one_eq_rightOverlappingLift}
identify the first two adjacent translates with the left and right overlapping
lifts of the same pair-indexed bond. Positivity makes this bond Hermitian, and
\leanid{MPOTensor.EtaLocalStructureData.overlappingLifts_pairBond_comm}
transports length-three commutation to the abstract lifts. The formal
hypotheses of the cited lemma therefore require no additional assumption.

The application to the pair-indexed bond is
\leanid{MPOTensor.EtaLocalStructureData.exists_unitary_blockActions_of_pairBond}.
It supplies one unitary middle-space identification and the two Hermitian
block actions required by the cited lemma.\footnote{The application theorem is
in \doclink{TNLean/MPS/MPDO/CommutingFormSpatialBridge}; its coordinate and
commutation hypotheses are supplied by
\doclink{TNLean/MPS/MPDO/CommutingOverlappingCoordinates}.}

\subsection{Local positive bond decomposition}

The local neighboring operators are extracted without assuming a Hayashi
decomposition. The theorem
\leanid{MPOTensor.EtaLocalStructureData.exists_positive_eta_pairBond_decomposition}
uses the same one-site unitary $U$ at both ends of the fixed bond and proves
\begin{align}
    (U^*\otimes U^*)B(U\otimes U)
        &=
        \bigoplus_{k,h}
        \Id_{H_{k,l}}\otimes\eta_{k,h}\otimes\Id_{H_{h,r}},
    \label{eq:cpsv16_commuting_form_to_sal_positive_eta_bond_decomposition}\\
    \eta_{k,h}
        &\succeq0.
    \label{eq:cpsv16_commuting_form_to_sal_positive_eta_blocks}
\end{align}
Because the bond and unitary are independent of the chain length, these
operators define one neighboring family for every translate.

The bond-local cyclic product is formalized for every $N\geq2$. The theorem
\leanid{MPOTensor.reindex_product_embedLocalOperator_of_etaPair_decomposition}
constructs the dependent cyclic regrouping
\begin{align}
    H^{\otimes N}
        &\simeq
        \bigoplus_{k_0,\ldots,k_{N-1}}
        \bigotimes_{n=0}^{N-1}
        \left(H_{k_n,r}\otimes H_{k_{n+1},l}\right),
    \qquad k_N=k_0,
    \label{eq:cpsv16_commuting_form_to_sal_cyclic_hilbert_regrouping}
\end{align}
and identifies the translated bond product as
\begin{align}
    \prod_{n=0}^{N-1}B_{n,n+1}
        &\cong
        \bigoplus_{k_0,\ldots,k_{N-1}}
        \bigotimes_{n=0}^{N-1}\eta_{k_n,k_{n+1}},
    \qquad k_N=k_0.
    \label{eq:cpsv16_commuting_form_to_sal_cyclic_eta_product}
\end{align}
The result includes $N=2$, where the two windows have the opposite cyclic
orders $(0,1)$ and $(1,0)$.

Tensor-power conjugation and the positive realization scalar are then combined
with the cyclic product. The theorem
\leanid{MPOTensor.EtaLocalStructureData.%
exists_positive_cyclic_eta_block_decomposition}
chooses the unitary, sector dimensions, and positive neighboring family once
and proves, for every $N\geq2$,
\begin{align}
    E_N^*(U^*)^{\otimes N}
        \sigma^{(N)}(\mathcal K)
        U^{\otimes N}E_N
        &=
        c_N
        \bigoplus_{k_0,\ldots,k_{N-1}}
        \bigotimes_{n=0}^{N-1}\eta_{k_n,k_{n+1}},
    \label{eq:cpsv16_commuting_form_to_sal_scaled_cyclic_eta_decomposition}\\
    c_N
        &>0,
    \qquad k_N=k_0.
    \label{eq:cpsv16_commuting_form_to_sal_scaled_cyclic_eta_conditions}
\end{align}

\subsection{Local correction to the proportionality scalar}

The converse argument at lines~1603--1605 of
Ref.~\cite{Cirac2016MPDO_arXiv} begins with proportionality between
$\sigma^{(N)}(\mathcal K)$ and the bond product, but then invokes the
coefficient-free source equation labelled \emph{sigmaNK2}. The formal theorem
instead proves the proportionality-preserving identity
(\ref{eq:cpsv16_commuting_form_to_sal_scaled_cyclic_eta_decomposition}). It
does not prove the coefficient-free equation printed at lines~1581--1588 of
Ref.~\cite{Cirac2016MPDO_arXiv}.

Proportional commuting-bond form alone cannot remove the coefficient. Take
physical dimension $d=1$, bond dimension $D=2$, and sole matrix-product letter
$M^{00}=I_2$. Then
\begin{align}
    \sigma^{(N)}
        &=\tr(I_2^N)
         =2.
    \label{eq:cpsv16_commuting_form_to_sal_scalar_counterexample_mpo}
\end{align}
The scalar bond $B=1$ realizes this family with $c_N=2$ for every $N$. Since
the physical space is one-dimensional, every positive neighboring
decomposition has one one-dimensional sector and one scalar $a\geq0$. A
coefficient-free identity at lengths two and four would require
\begin{align}
    a^2
        &=2,
    \label{eq:cpsv16_commuting_form_to_sal_scalar_counterexample_length_two}\\
    a^4
        &=2.
    \label{eq:cpsv16_commuting_form_to_sal_scalar_counterexample_length_four}
\end{align}
These equalities are incompatible. Thus a fixed local rescaling does not
follow from the displayed proportionality data.

This example is not a normal injective block and therefore does not contradict
the source setting. It shows that a scalar-rescaling proof would need normal
injectivity. Such a proof would have to compare the matrix-product family with
the fixed bond-product family and establish
\begin{align}
    c_N
        &=a^N
        \qquad\text{for one }a>0.
    \label{eq:cpsv16_commuting_form_to_sal_exponential_realization_scalar}
\end{align}
Only then could $a$ be absorbed once into the bond, or equivalently into the
neighboring operators, independently of $N$. The completed proof takes a
different route: it directly compares normalized marginals of the selected
fixed-product tensor with those of the source tensor.

\subsection{The selected fixed-product tensor}

The tensor \leanid{MPOTensor.cyclicEdgeWeightTensor} has virtual dimension
$d^2$ and generates one prescribed cyclic nearest-neighbor scalar weight at
every positive length. The datum
\leanid{MPOTensor.TranslationInvariantBondData.FixedProductTensorData}
records one positive-dimensional tensor whose operator family equals the
product of translates of the fixed bond for every $N\geq2$. Given this datum,
\leanid{MPOTensor.EtaLocalStructureData.%
eventuallyNonzeroProportionalMPV₂_fixedProductTensor}
derives eventual nonzero proportionality to the source tensor without a
length-one hypothesis.

The theorem
\leanid{MPOTensor.TranslationInvariantBondData.fixedProductTensorData}
constructs the fixed tensor directly from arbitrary
\leanid{MPOTensor.TranslationInvariantBondData}. It assembles the
chain-independent sector coordinates and neighboring operators into a cyclic
scalar weight in transformed one-site coordinates, then conjugates the tensor
back through the same one-site unitary. The construction uses only positivity
and commutation of the periodic translates.

The selected tensor retains the positive physical-sector factorization used
in its construction. Specifically,
\leanid{MPOTensor.TranslationInvariantBondData.%
exists_positive_physicalSectorFactorization_fixedProductTensorData}
factorizes the exact tensor in
\leanid{MPOTensor.TranslationInvariantBondData.fixedProductTensorData}.
Its neighboring operators are positive semidefinite, and its reconstructed
physical bond is the original commuting bond. No equality of closed
matrix-product operators is used to replace one local tensor by another.

This construction supplies only a fixed representative. It neither factorizes
the original source tensor nor proves a normality or blocking comparison
between that tensor and the selected fixed tensor.

The coordinate qualification is essential. In physical dimension $d=2$, let
$X$ be the Pauli flip and set $B=\Id+X\otimes X$. Then $B\succeq0$, all
periodic translates commute, and
\begin{align}
    \bra{0^N}
        \left(\prod_n B_{n,n+1}\right)
        \ket{0^N}
        &=2
        \qquad(N\geq2).
    \label{eq:cpsv16_commuting_form_to_sal_pauli_bond_zero_boundary_entry}
\end{align}
At $N=2$, the repository's periodic product is
\begin{align}
    B^2
        &=2B.
    \label{eq:cpsv16_commuting_form_to_sal_pauli_bond_length_two_product}
\end{align}
For $N\geq3$, the empty edge set and the full cycle give the two terms with
zero boundary. A scalar edge weight in the original coordinates would make
the entry $a^N$, so lengths two and three would require
\begin{align}
    a^2
        &=2,
    \label{eq:cpsv16_commuting_form_to_sal_pauli_scalar_length_two}\\
    a^3
        &=2.
    \label{eq:cpsv16_commuting_form_to_sal_pauli_scalar_length_three}
\end{align}
These equalities are incompatible. An arbitrary fixed bond therefore need not
satisfy
\leanid{MPOTensor.TranslationInvariantBondData.CyclicEdgeWeightForm} in the
original coordinates. The source equation labelled \emph{sigmaNK2} in
Ref.~\cite{Cirac2016MPDO_arXiv} gives the scalar form only after the
one-site unitary coordinate change.

An arbitrary tensor representing the bond products need not become normal
after one positive rescaling. Adjoining a zero virtual direct-sum block changes
neither the closed operator at any positive length nor the product identity,
but no nonzero rescaling removes that invariant summand. Proposition
\emph{3to4} of Ref.~\cite{Cirac2016MPDO_arXiv} does not assert normality or
uniqueness for an arbitrary representation.

No canonical-form reduction is needed for the bond selected in the SAL proof.
The source physical-sector factorization gives the coefficient-one identity
directly, so the original source tensor represents the bond product. This is
\leanid{MPOTensor.PhysicalSectorFactorization.fixedProductTensorData}. Under
the normal single-block hypothesis, normality and the exact all-length identity
are recorded by
\leanid{MPOTensor.exists_normal_fixedProductTensorData_of_isSAL}. The
source-selected presentation therefore has rescaling $s=1$.

For a different proportional commuting-bond presentation, the conditional
theorem
\leanid{MPOTensor.EtaLocalStructureData.%
mpo_eq_positive_power_product_of_normal_smul}
shows that if $sC$ is normal for one $s>0$, then
\begin{align}
    \sigma^{(N)}(\mathcal K)
        &=s^N\prod_n B_{n,n+1}
        \qquad(N\geq2).
    \label{eq:cpsv16_commuting_form_to_sal_normal_rescaling_product_identity}
\end{align}
The normal-tensor fundamental theorem first permits a unit phase $\zeta^N$.
The self-overlap of the normal source tends to one, so the fixed product is
nonzero at two sufficiently large consecutive lengths. Positive realization
constants at those lengths force $\zeta=1$. This strengthening applies when a
normal rescaling is already known. Appendix~C.2 of
Ref.~\cite{Cirac2016MPDO_arXiv} neither claims nor needs the existence of such
a rescaling for every proportional commuting-bond presentation.

The scaled identity uses only the fixed commuting bond and its positive
realization scalar. It does not assert zero correlation length, rank-one trace
factorization, a quantum Markov decomposition, saturation of the area law, or
Proposition \emph{4to2} of Ref.~\cite{Cirac2016MPDO_arXiv}.

\subsection{The fourth-region calculation}

Tracing one physical site from the length-$N$ cyclic eta product leaves the
coupled sum
\begin{align}
    \sum_q
        \tr_{H_{q,l}}\!\left(\eta_{k_{N-1},q}\right)
        \otimes
        \tr_{H_{q,r}}\!\left(\eta_{q,k_1}\right).
    \label{eq:cpsv16_commuting_form_to_sal_one_site_traced_eta_boundary_sum}
\end{align}
The same sector $q$ occurs in both factors, so rank-one factorization does not
separate the expression into independent sums. At line~1613,
Ref.~\cite{Cirac2016MPDO_arXiv} instead invokes ZCL to replace this marginal
by a trace of two boundary sites on the length-$(N+1)$ chain. In edge
coordinates, the middle operator $\eta_{q,h}$ is then fully traced, while the
adjacent operators retain their right and left partial traces. The source
formula at line~1616 suppresses these partial-trace symbols.

The conditional algebraic identity after this replacement is
\leanid{Matrix.eta_fourth_region_trace_formula}. It assumes
\begin{align}
    \tr(\eta_{q,h})
        &=a_qb_h
    \label{eq:cpsv16_commuting_form_to_sal_rank_one_eta_trace_assumption}
\end{align}
and separates the two boundary sums. It does not prove the ZCL marginal
replacement, derive rank-one factorization, construct a quantum Markov
decomposition, or prove SAL.

The ZCL replacement is formalized separately. Let $E$ be the physical-trace
transfer and suppose
\begin{align}
    E^2
        &=\lambda E,
    \qquad \lambda>0.
    \label{eq:cpsv16_commuting_form_to_sal_source_zcl_transfer_identity}
\end{align}
For every $N\geq0$, the normalized $N$-site marginal obtained by tracing two
sites from the length-$(N+2)$ chain then equals the marginal obtained by
tracing one site from the length-$(N+1)$ chain. Entrywise,
\begin{align}
    \frac{\tr(M^{u,v}E^2)}
         {\tr(E^{N+2})}
        &=
        \frac{\lambda\tr(M^{u,v}E)}
             {\lambda\tr(E^{N+1})},
    \label{eq:cpsv16_commuting_form_to_sal_zcl_two_site_reduction}\\
        &=
        \frac{\tr(M^{u,v}E)}
             {\tr(E^{N+1})}.
    \label{eq:cpsv16_commuting_form_to_sal_zcl_one_site_reduction}
\end{align}
This is
\leanid{MPOTensor.reducedBlockState_succ_succ_eq_succ_of_isSourceZCL}.
It proves the first sentence of the argument at source line~1613, but it does
not separate the two boundary-sector sums.

A separate obstruction occurs at source line~1613. The phrase ``arguing as in
Lemmas \emph{propSN} and \emph{SALZCL}'' in
Ref.~\cite{Cirac2016MPDO_arXiv} does not show that source ZCL alone implies
(\ref{eq:cpsv16_commuting_form_to_sal_rank_one_eta_trace_assumption}). Lemma
\emph{propSN} derives primitivity from SAL, which is the conclusion sought in
Proposition \emph{4to2}. Lemma \emph{SALZCL} then uses the invalid
trace-power-to-rank-one inference recorded in
\path{docs/paper-gaps/cpgsv17_pf_rank_one.tex}.

Formally, source ZCL yields an idempotent normalized pairing operator. For the
opposite product, however, it yields only a square--cube identity and constant
positive-power traces. The four-sector example
\leanid{MPOTensor.ActiveSectorSpanningCounterexample.%
virtual_spanning_does_not_force_rectangular_remainder_zero}
satisfies the recorded positivity, primitivity, trace-one, and spanning
conditions while its trace matrix is not idempotent.

The capstone theorem
\leanid{MPOTensor.ActiveSectorSpanningCounterexample.%
tensor_has_sal_sourceZCL_and_non_rankOne_activeTraceMatrix}
adds SAL, TNLean's up-to-scalar source-ZCL predicate, and the source-selected
inverse-map factorization, and proves that the active trace matrix is not rank
one. This failure belongs to the selected four-sector description. The same
tensor has a one-sector physical factorization with a positive trace-one
neighboring operator, formalized by
\leanid{MPOTensor.ActiveSectorSpanningCounterexample.%
exists_oneSector_neighboringTraceFactorization}.
Consequently, the four-sector trace matrix does not obstruct the existence of
another factorization.

The raw tensor in this example is not the normal tensor assumed in
Appendix~C.2, Case~I, of Ref.~\cite{Cirac2016MPDO_arXiv}; its normal
representative also fails the paper's literal ZCL diagram. The example
therefore marks two proof-route boundaries rather than refuting Lemma~C.5
under its full standing hypotheses. The claimed derivation of rank one from
source ZCL alone remains unavailable. The completed proof does not use that
derivation and instead obtains a positive rank-one factorization of the
two-step coefficient.

\subsection{Cyclic-active sectors and the two-step coefficient}

Primitivity cannot be assumed for an arbitrary selected spatial
factorization. The theorem
\leanid{MPOTensor.CommutingBondTraceMatrixObstruction.%
normal_sourceZCL_fixed_commutingBond_does_not_force_traceMatrix_primitive}
gives an injective, normal, bond-dimension-one product MPDO with source ZCL and
one exact fixed positive commuting-bond presentation. Its selected two-sector
factorization has positive neighboring operators and trace matrix
\begin{align}
    T
        &=
        \begin{pmatrix}
            1&0\\
            0&0
        \end{pmatrix}.
    \label{eq:cpsv16_commuting_form_to_sal_nonprimitive_selected_trace_matrix}
\end{align}
The second physical summand is inactive but retained by the nonminimal
middle-space decomposition, so the full matrix is not primitive. The same
product tensor has a one-sector factorization with primitive trace matrix
\begin{align}
    T_{\mathrm{min}}
        &=
        \begin{pmatrix}
            1
        \end{pmatrix}.
    \label{eq:cpsv16_commuting_form_to_sal_primitive_minimal_trace_matrix}
\end{align}
This example does not exclude a primitive spatial factorization. It shows that
the direct-sum tensor-factor decomposition must be supplemented by a
source-faithful minimal or visible-sector selection before primitivity can be
used.

Such a selection is available without SAL. A sector is cyclic-active if it has
an outgoing nonzero neighboring operator followed by a possibly empty return
path. The explicit first edge excludes the empty reflexive walk. This is
\leanid{MPOTensor.PhysicalSectorFactorization.IsCyclicActiveSector}. Every
nonzero finite cyclic neighboring product uses only cyclic-active sectors.
Deleting all other sectors therefore preserves every finite cyclic bond
product; see
\leanid{MPOTensor.PhysicalSectorFactorization.%
cyclicNeighboringProduct_eq_zero_of_not_isCyclicActiveSector}.

For an injective tensor with positive semidefinite neighboring operators, the
cyclic-active sector set is nonempty and its trace matrix is primitive. The
proof uses nondegeneracy of the matrix trace pairing twice. First, every
nonzero virtual sector lies on a directed two-cycle. All virtual matrices
outside the cyclic-active set therefore vanish, and the retained virtual
matrices still span the full matrix algebra. A directed-cut argument then
gives one strongly connected component. Second, each retained sector has a
return of length two, and each retained edge closes through a third retained
sector. These coprime return lengths give aperiodicity. This result is
\leanid{MPOTensor.PhysicalSectorFactorization.%
cyclicActiveSectorTraceMatrix_isPrimitive}. It uses no inverse-map sector
weights, SAL assumption, or Hayashi decomposition.

Source ZCL gives a positive normalization
$\widehat T=\lambda^{-1}T$ of the restricted trace matrix satisfying
\begin{align}
    \widehat T^2
        &=\widehat T^3,
    \label{eq:cpsv16_commuting_form_to_sal_normalized_trace_square_cube}\\
    \operatorname{rank}(\widehat T^2)
        &=1,
    \label{eq:cpsv16_commuting_form_to_sal_normalized_trace_square_rank}\\
    \tr(\widehat T^m)
        &=\tr(\widehat T)
        \qquad(m\geq1).
    \label{eq:cpsv16_commuting_form_to_sal_normalized_trace_power_traces}
\end{align}
The combined result is
\leanid{MPOTensor.PhysicalSectorFactorization.%
exists_normalized_cyclicActiveSectorTraceMatrix_rank_pow_two_eq_one_of_isSourceZCL}.
It concerns only the restricted matrix. It neither makes the unreduced trace
matrix primitive nor asserts that $\widehat T$ has rank one.

This distinction determines how the result enters Proposition \emph{4to2} of
Ref.~\cite{Cirac2016MPDO_arXiv}. The two-boundary expression at source
lines~1613--1616 contains the one-step coefficient
\begin{align}
    T_{q,h}
        &=\tr(\eta_{q,h}),
    \label{eq:cpsv16_commuting_form_to_sal_one_step_trace_coefficient}
\end{align}
not $T^2$. Applying the ZCL marginal replacement once more traces three
boundary sites of a chain longer by two sites and produces two successive
fully traced middle edges:
\begin{align}
    (T^2)_{q,h}
        &=
        \sum_{r\in C}
        \tr(\eta_{q,r})\tr(\eta_{r,h}).
    \label{eq:cpsv16_commuting_form_to_sal_two_step_trace_coefficient}
\end{align}

The second iteration of the normalized marginal identity is
\leanid{MPOTensor.reducedBlockState_add_three_eq_succ_of_isSourceZCL}.
The coefficient calculation in cyclic-active physical-sector coordinates is
\leanid{MPOTensor.PhysicalSectorFactorization.%
sum_cyclicActive_partialTraceRight_threeSiteSectorClosure}.
Its normalized version,
\leanid{MPOTensor.PhysicalSectorFactorization.%
sum_cyclicActive_normalized_trace_mul_trace_eq_pow_two},
scales each traced edge by $\lambda^{-1}$ and yields
$(\widehat T^2)_{q,h}$. Moreover,
\leanid{MPOTensor.PhysicalSectorFactorization.%
map_cyclicNeighboringProduct_eq_zero_of_not_forall_isCyclicActiveSector}
shows that arbitrary linear coordinate contraction preserves the deletion of
every vanishing cyclic summand. The rank-one statement in
(\ref{eq:cpsv16_commuting_form_to_sal_normalized_trace_square_rank}) therefore
supplies the separating boundary coefficient in the three-boundary formula.

The complete physical-sector direct sum now carries the normalized marginal
equality while retaining the two outer partial traces. This is
\leanid{MPOTensor.PhysicalSectorFactorization.%
exists_cyclicActiveFourthRegion_formula_of_isSourceZCL}.
Separating the retained chain at an arbitrary distinguished site, normalizing
the positive left and right path factors, and collecting their trace products
gives
\leanid{MPOTensor.PhysicalSectorFactorization.%
exists_hayashiMarkovDecomposition_cyclicActiveCut_of_isSourceZCL}.
Thus the Markov decomposition is available in physical-sector coordinates for
arbitrary left and right lengths.

Substituting the admissible entropy-cut lengths, applying the all-cut Markov
criterion, and transporting through the physical isometry gives
\leanid{MPOTensor.PhysicalSectorFactorization.isSAL_of_isSourceZCL}.
The cut calculation and basis transport are complete under an explicit
positive physical-sector factorization. The older conditional theorem
\leanid{Matrix.eta_fourth_region_trace_formula} remains a coordinate model for
a one-step rank-one coefficient; it is not used to identify $T_C$ with
$T_C^2$.

\subsection{Transfer from the selected representative}

The explicit-factorization theorem
\leanid{MPOTensor.PhysicalSectorFactorization.isSAL_of_isSourceZCL} is
scope-restricted when considered alone: it applies to the tensor that carries
the factorization and therefore does not by itself prove
\begin{align}
    (\text{one translation-invariant commuting bond})
        +(\text{source ZCL})
        &\Longrightarrow
        \text{SAL}.
    \label{eq:cpsv16_commuting_form_to_sal_source_implication}
\end{align}
For the source theorem, formalized by
\leanid{MPOTensor.EtaLocalStructureData.isSAL_of_isSourceZCL}, the required
positive physical-sector factorization is constructed on the selected fixed
tensor obtained from the bond in the source equation labelled
\emph{sigmaNK2}, at lines~1581--1605 of
Ref.~\cite{Cirac2016MPDO_arXiv}. The remaining comparison transfers normalized
marginals rather than assuming injectivity, normality, or ZCL for the selected
tensor.

The recurrent-sector comparison requires no normality or blocking comparison
with the selected representative. Express the original injective tensor in
the physical coordinates of the selected factorization. The length-two
closed-operator identity eliminates all cross-sector matrices and all matrices
outside the positive-length cyclic support. The remaining matrices still span
the full auxiliary matrix algebra.

Positivity gives a nonzero closed coefficient through every cyclic-active
sector. The length-three identity sends every zero neighboring operator to a
zero ordered product of the corresponding two sector spaces. A proper
reachability cut would then split the full matrix algebra into two nonzero
subspaces whose products vanish in one order, which is impossible. The theorem
\leanid{MPOTensor.EtaLocalStructureData.%
selectedFixedProduct_cyclicActive_reflTransGen_neighboringOperator_ne_zero}
therefore proves that the cyclic-active sectors of the selected factorization
form one recurrent component. It does not assert injectivity, normality, or a
virtual gauge for the selected fixed tensor. This closes the selected-sector
visibility step tracked in issue~\#5128.

The coefficient comparison at source lines~1606--1617 of
Ref.~\cite{Cirac2016MPDO_arXiv} is also formalized. Expanding the normalized
reduced state after tracing one suffix site and after tracing two suffix sites
produces surviving boundary contractions with coefficients $T_C^2$ and
$T_C^3$, respectively. Recurrence shows that these powers have the same
positive support. Irreducibility then makes every entry of $T_C^2$ strictly
positive. A positive two-edge summand supplies a retained word of common
length three for every ordered pair, and comparison at that fixed length gives
\begin{align}
    \widehat T^2
        &=\widehat T^3,
    \qquad
    \widehat T=\lambda^{-1}T_C.
    \label{eq:cpsv16_commuting_form_to_sal_selected_trace_square_cube}
\end{align}
Here the common normalization is the ratio of the selected length-five and
length-four traces. This result is
\leanid{MPOTensor.EtaLocalStructureData.%
exists_normalized_selectedFixedProduct_cyclicActiveSectorTraceMatrix_pow_two_eq_pow_three}.
It uses source ZCL and injectivity of the original tensor, not of the selected
fixed-product tensor. This closes the adjacent coefficient-extraction step
tracked in issue~\#5129.

The square--cube identity combines with the strictly positive two-step
coefficient. Nonemptiness of the cyclic-active sector set follows from the
spanning family of original one-site matrices: if the sector set were empty,
that family would have zero span instead of the full auxiliary matrix algebra.
The normalized trace matrix is therefore primitive, with its second power
strictly positive. The theorem
\leanid{Matrix.IsPrimitive.%
exists_pos_pow_two_eq_vecMulVec_of_pow_two_eq_pow_three}
then gives strictly positive vectors $a$ and $b$ such that
\begin{align}
    (\lambda^{-1}T_C)^2
        &=ab^{\mathsf T},
    \qquad a>0,\quad b>0.
    \label{eq:cpsv16_commuting_form_to_sal_positive_two_step_rank_one}
\end{align}

The one-step and iterated marginal identities, together with positivity of the
selected traces, are transported separately by
\leanid{MPOTensor.EtaLocalStructureData.%
selectedFixedProduct_reducedBlockState_succ_succ_eq_succ_of_isSourceZCL},
\leanid{MPOTensor.EtaLocalStructureData.%
selectedFixedProduct_reducedBlockState_add_three_eq_succ_of_isSourceZCL},
and
\leanid{MPOTensor.EtaLocalStructureData.%
exists_pos_trace_mpo_selectedFixedProduct_of_isSourceZCL}.

Substituting
(\ref{eq:cpsv16_commuting_form_to_sal_positive_two_step_rank_one}) into the
three-suffix formula and regrouping at an arbitrary distinguished site gives
the Hayashi direct sum. Positive proportionality identifies the selected
normalized marginal with the corresponding marginal of the original source
tensor in the same one-site physical coordinates. This is
\leanid{MPOTensor.EtaLocalStructureData.%
exists_hayashiMarkovDecomposition_selectedPhysicalCut_of_isSourceZCL}.
It assumes injectivity of the source tensor, an eta-local structure, and source
ZCL. It assumes neither MPDO positivity at this intermediate stage nor
injectivity, normality, or zero correlation length for the selected
fixed-product tensor. MPDO positivity enters only in the final SAL theorem,
where it supplies positivity of the source operators in the selected physical
coordinates.

Applying the all-cut Markov criterion to these physical-coordinate marginals
and transporting saturation through the one-site isometry gives
\leanid{MPOTensor.EtaLocalStructureData.isSAL_of_isSourceZCL}. Positivity of
the physical-coordinate periodic operators follows from the MPDO hypothesis
by isometric conjugation. Their traces equal the original traces, which source
ZCL makes nonzero at every positive length. The theorem assumes precisely the
source tensor, its MPDO and injectivity hypotheses, the fixed eta-local
structure, and source ZCL. It assumes no injectivity, normality, ZCL, or SAL
property of the selected fixed-product tensor.

\subsection{Scope restriction for sectorwise source ZCL}

The proposition labelled \emph{prop4to2}, at lines~1801--1808 of
Ref.~\cite{Cirac2016MPDO_arXiv}, states that GSNNCH form alone implies SAL.
Its proof applies Proposition \emph{4to2} to each summand, although
Proposition \emph{4to2} assumes ZCL at source lines~1597--1600. The surrounding
main theorem includes ZCL in condition~(iii), at source lines~851--893. The
printed proposition therefore omits a hypothesis used in its proof.

The final sentence of that proof is a separate valid claim: a finite
orthogonal direct sum on fixed one-site subspaces satisfies SAL if every
summand satisfies SAL. The formal scope-restricted theorem assumes SAL for
the summands, a positive proportional local orthogonal-sum identity at every
positive chain length, supplied one-site projections supporting every
nonempty sector marginal, and positive natural multiplicities. It must not be
presented as a proof that GSNNCH alone implies SAL.

The structure
\leanid{MPOTensor.ProportionalOrthogonalCommutingSectorFamily}
records fixed one-site projections and commuting sector products whose
periodic products realize the sector operators up to positive scalars.
Projection idempotence, bond support, and positive proportional realization
show that the normalized sector chain and every nonempty sector marginal on a
chain of length at least two are supported on the matching first-site
projection. These results are
\leanid{MPOTensor.ProportionalOrthogonalCommutingSectorFamily.%
firstSiteMatrix_mul_normalizedMPO}
and
\leanid{MPOTensor.ProportionalOrthogonalCommutingSectorFamily.%
firstSiteMatrix_mul_reducedBlockState}.
The scalar passes through the first-site action and cancels under trace
normalization. The earlier
\leanid{MPOTensor.OrthogonalCommutingSectorFamily} is the scalar-one
specialization. The chain-length restriction suffices because the
mutual-information comparison in SAL uses only chains of length at least four.

The sector-family structure contains neither the outer tensor nor its
multiplicities, so it cannot itself provide the proportional outer sum at
every positive length. The theorem
\leanid{MPOTensor.isSAL_of_proportionalLocalOrthogonalSum} assumes positive
numbers $c_N$ satisfying
\begin{align}
    \sigma^{(N)}(M)
        &=c_N\sum_s n_s\sigma^{(N)}(K_s),
    \qquad c_N>0,
    \label{eq:cpsv16_commuting_form_to_sal_proportional_orthogonal_outer_sum}
\end{align}
and cancels $c_N$ from the normalized chain and all its marginals. It assumes
sectorwise SAL and proves the noncircular local direct-sum implication. It
neither derives sectorwise SAL nor uses ZCL. The exact theorem
\leanid{MPOTensor.isSAL_of_localOrthogonalSum} is the specialization $c_N=1$.

Ambient source ZCL supplies the missing sectorwise source-ZCL hypothesis in
the canonical BNT specialization. For a simple tensor already in the selected
canonical form,
\leanid{MPOTensor.IsSimpleCanonicalForm.%
exists_commonWeightAbsorbedBasisMPOTensor_isSourceZCL}
chooses the canonical BNT data, proves that the weights over each normal
representative are independent of the copy index, and proves source ZCL for
every absorbed representative. Its final step uses
\leanid{MPOTensor.commonWeightAbsorbedBasisMPOTensor_isSourceZCL}.
Thus every canonical absorbed simple-BNT sector has source ZCL under ambient
source ZCL. This is exactly the sectorwise premise used when Proposition
\emph{4to2} is applied in the proof of Proposition \emph{prop4to2} in
Ref.~\cite{Cirac2016MPDO_arXiv}.

This conclusion applies to the canonical absorbed presentation, not to an
arbitrary GSNNCH presentation of the same finite-chain operators. Source ZCL
is a property of the tensor's physical-trace transfer, whereas GSNNCH records
only the represented periodic operators. For example, in physical dimension
one, the tensor $(1)$ and the virtual direct sum
\begin{align}
    (1)\oplus
    \begin{pmatrix}
        0&1\\
        0&0
    \end{pmatrix}
    \label{eq:cpsv16_commuting_form_to_sal_nilpotent_virtual_extension}
\end{align}
generate the same scalar operator at every positive length. The first tensor
has source ZCL, while the second transfer contains a nonzero nilpotent summand
and satisfies no identity
\begin{align}
    E^2
        &=\lambda E
        \qquad\text{with }\lambda>0.
    \label{eq:cpsv16_commuting_form_to_sal_failed_zcl_for_nilpotent_extension}
\end{align}
Tensor-level source ZCL therefore cannot be inferred from GSNNCH alone without
selecting the canonical absorbed sectors.

The downstream theorem
\leanid{MPOTensor.%
isSAL_of_proportionalOrthogonalCommutingSectorFamily_of_sectorwise_isSourceZCL}
requires only that each sector generate MPDOs, be one-site injective, and have
source ZCL. It does not require normality. Source ZCL forces nonzero bond
dimension through \leanid{MPOTensor.IsSourceZCL.bondDim_ne_zero}, so no
normality assumption is needed to obtain the nonzero-dimensional bond-space
hypothesis in the all-cut Markov argument.

For canonical absorbed sectors, the standing biCF one-letter span and
copy-independent weights give one-site injectivity through
\leanid{MPOTensor.commonWeightAbsorbedBasisMPOTensor_isInjective}. Ambient
source ZCL descends to each absorbed sector through
\leanid{MPOTensor.commonWeightAbsorbedBasisMPOTensor_isSourceZCL}. The only
remaining positivity question is the chain length not covered by the
commuting-bond product formula.

The proportional orthogonal commuting sector family records a positive
commuting bond whose periodic product is proportional to the sector operator
by a positive scalar for every $N\geq2$. Hence
\leanid{MPOTensor.ProportionalOrthogonalCommutingSectorFamily.%
mpo_posSemidef_of_two_le}
proves positivity at every such length. The theorem
\leanid{MPOTensor.ProportionalOrthogonalCommutingSectorFamily.%
isMPDO_of_mpo_one_pos}
shows that the remaining sharp interface is
\begin{align}
    \rho^{(1)}(K_s)
        &\succeq0.
    \label{eq:cpsv16_commuting_form_to_sal_absorbed_sector_length_one_positivity}
\end{align}
Together, these facts give full sector MPDO positivity. No positivity for
$N\geq2$ and no spectral normality of an absorbed sector is assumed
separately.

The corrected canonical statement is
\leanid{MPOTensor.%
isSAL_of_commonWeightAbsorbedBasisMPOTensor_of_proportionalSectors_of_isSourceZCL}.
It assumes the selected canonical decomposition, the reconstructing
block-diagonal gauge and positive-length MPV equality, copy-independent
weights, nonnilpotence of the representative physical-trace transfers, the
standing biCF one-letter product-algebra span, the proportional orthogonal
commuting sector family, one-site positivity of every absorbed sector, and
ambient source ZCL. The exact canonical BNT outer sum, one-site injectivity,
sectorwise source ZCL, and the preceding positivity theorem then supply all
hypotheses of the general orthogonal-sector SAL theorem. The conclusion is SAL
at every original cut.

The length-one premise cannot be removed from the source commuting-product
display alone: TNLean's unambiguous periodic bond realization begins at
$N=2$. The standing Case~II biCF hypothesis already supplies the one-letter
simultaneous span required for injectivity at every original cut.

\section{Verdict}

Proposition \emph{4to2} of Ref.~\cite{Cirac2016MPDO_arXiv} is formalized by
\leanid{MPOTensor.EtaLocalStructureData.isSAL_of_isSourceZCL}. The positive
fixed-product representative, its cyclic-active recurrent component, the
normalized square--cube identity, and the positive rank-one factorization of
the two-step coefficient are derived from the source tensor. The resulting
Hayashi decomposition holds at every one-site cut, and isometric invariance
returns SAL to the original physical coordinates. Normality is not a
hypothesis of this theorem.

For Proposition \emph{prop4to2} of Ref.~\cite{Cirac2016MPDO_arXiv}, the
positive proportional outer-sum theorem reduces SAL to sectorwise SAL. In the
canonical Case~II application, ambient source ZCL supplies sectorwise source
ZCL, the biCF one-letter span supplies sector injectivity, and one-site
positivity completes sector positivity because the commuting bonds already
give positivity for all $N\geq2$. The corrected canonical theorem is complete
at the ambient-source-ZCL and one-site-positivity boundary. It requires no
independent normality or full MPDO premise for the absorbed sectors.

The literal Proposition \emph{prop4to2} remains unproved as printed. Its
statement assumes only GSNNCH, whereas Proposition \emph{4to2}, invoked in its
proof, also assumes ZCL. Condition~(iii) of the surrounding main theorem
supplies ambient ZCL, and the corrected ZCL-conditioned canonical-BNT
application is covered by the preceding argument. An arbitrary GSNNCH
presentation does not inherit tensor-level source ZCL, as the nilpotent-summand
example
(\ref{eq:cpsv16_commuting_form_to_sal_nilpotent_virtual_extension}) shows.
Ambient source ZCL therefore remains a separate scope restriction on the
available implication.

\section{Completed proof}

\begin{enumerate}
    \item Apply the all-cut Markov criterion to the decompositions supplied by
    \leanid{MPOTensor.EtaLocalStructureData.%
    exists_hayashiMarkovDecomposition_selectedPhysicalCut_of_isSourceZCL}.
    Positivity and nonvanishing of the physical-coordinate periodic operators
    follow from the MPDO and source-ZCL hypotheses.

    \item Transport SAL through the selected one-site physical isometry back
    to the original coordinates. These two steps are formalized by
    \leanid{MPOTensor.EtaLocalStructureData.isSAL_of_isSourceZCL}, whose
    blueprint entry carries a \texttt{\string\leanok} marker.
\end{enumerate}

\bibliographystyle{alpha}
\bibliography{references}

\end{document}
